\documentclass[lualatex, paper={90mm, 160mm}]{jlreq}
\usepackage{tikz}
\usepackage{lltjext}
\usepackage{xcolor}
\usetikzlibrary{calc}

% カラーパレットの定義
\definecolor{DeepRed}{HTML}{660000}
\definecolor{RichGold}{HTML}{B8860B}
\definecolor{DarkBlue}{HTML}{001A33}

\begin{document}
\pagestyle{empty} 

\begin{tikzpicture}[remember picture, overlay]
  % 1. 背景全体を濃紺で塗りつぶす
  \fill[DarkBlue] (current page.south west) rectangle (current page.north east);

  % 2. 背景の曲線模様（雲気・流水をモチーフにしたサイン波とベジェ曲線）
  \begin{scope}[opacity=0.8, line width=3pt]
    % 深紅のサイン波
    \foreach \i in {1,...,12} {
      \draw[DeepRed] 
        ($(current page.south west) + (0, \i*15mm - 30mm)$) 
        sin ($(current page.south) + (-15mm, \i*15mm + 10mm)$) 
        cos ($(current page.south) + (15mm, \i*15mm - 10mm)$) 
        sin ($(current page.south east) + (0, \i*15mm)$);
    }
    % 金色の流線（ベジェ曲線）
    \foreach \i in {1,...,6} {
      \draw[RichGold, line width=2pt] 
        ($(current page.north west) + (-10mm, -\i*25mm)$) 
        .. controls ($(current page.center) + (20mm, 30mm)$) and ($(current page.center) + (-20mm, -30mm)$) .. 
        ($(current page.south east) + (10mm, \i*25mm)$);
    }
  \end{scope}

  % 3. 表紙の二重枠線（金色に変更）
  \draw[line width=1.5pt, color=RichGold] ($(current page.south west) + (6mm, 6mm)$) rectangle ($(current page.north east) - (6mm, 6mm)$);
  \draw[line width=0.5pt, color=RichGold] ($(current page.south west) + (7.5mm, 7.5mm)$) rectangle ($(current page.north east) - (7.5mm, 7.5mm)$);

  % 4. タイトル（縦書き） - 視認性確保のための白半透明背景
  \node[anchor=north, fill=white, fill opacity=0.85, text opacity=1, inner sep=5mm, rounded corners=2mm] at ($(current page.center) + (15mm, 50mm)$) {
    \begin{minipage}<t>{15em}
      \begin{center}
        {\fontsize{28pt}{40pt}\selectfont \textbf{走れメロス}\par}
      \end{center}
    \end{minipage}
  };

  % 5. 作者名（縦書き）
  \node[anchor=north, fill=white, fill opacity=0.85, text opacity=1, inner sep=3mm, rounded corners=2mm] at ($(current page.center) + (-15mm, -10mm)$) {
    \begin{minipage}<t>{10em}
      \begin{center}
        {\fontsize{16pt}{24pt}\selectfont \textbf{太宰 治}\par}
      \end{center}
    \end{minipage}
  };
\end{tikzpicture}

\end{document}